\newcommand{\DataDate}{2026-10-14}
\newcommand{\FreezeDate}{2026-10-14}
\newcommand{\FirstMonth}{2006-02}
\newcommand{\LastMonth}{2026-09}
\newcommand{\NMonths}{248}
\newcommand{\NFormations}{249}
\newcommand{\AvgUniverse}{31.6}
\newcommand{\NTickers}{58}
\newcommand{\NMembersEver}{65}
\newcommand{\NMissing}{7}
\newcommand{\CoverageMean}{96.9\%}
\newcommand{\NChanges}{43}
\newcommand{\CostBps}{10}
\newcommand{\BorrowBps}{25}
\newcommand{\NTrials}{36}
\newcommand{\SRStar}{0.43}
\newcommand{\DaxCagr}{7.4\%}
\newcommand{\DaxSharpe}{0.43}
\newcommand{\DaxMaxDD}{$-$52.4\%}
\newcommand{\SpxCagr}{11.5\%}
\newcommand{\SpxSharpe}{0.74}
\newcommand{\CompCagr}{9.8\%}
\newcommand{\CompSharpe}{0.67}
\newcommand{\SurvBias}{3.8\%}
\newcommand{\NBigMoves}{23}
\newcommand{\NFixes}{9}
\newcommand{\FindBest}{The strongest signal is \emph{Momentum 12-1}, whose long-short portfolio earns 7.2\% per year after costs (Sharpe 0.40); its mean return is only marginally significant (10\% level).}
\newcommand{\FindWorst}{The weakest is \emph{Short-term reversal} with a net Sharpe ratio of $-$0.29; it loses money mostly through trading costs, since its gross spread is only $-$0.4\% per year at a turnover of 16.3 times the portfolio per year.}
\newcommand{\FindHolm}{After a Holm correction for testing six strategies, no long-short premium remains significant at the 5\% level.}
\newcommand{\FindDSR}{Across all 36 specifications that were run, an unskilled researcher would expect a best Sharpe ratio of 0.43 by luck alone. No strategy has a deflated Sharpe ratio above 0.95, so none of them can be called a discovery on this data.}
\newcommand{\FindBeta}{Five of six long-short portfolios carry a negative market beta (lowest: Low beta at $-$0.76), because the bottom tercile of these signals is dominated by volatile, high-beta stocks. In a rising market that drag hides part of the premium: after controlling for the DAX, the CAPM alpha is significant at 5\% for Momentum 12-1, Low beta and 52-week high.}
\newcommand{\FindComposite}{The long-only composite portfolio returns 9.8\% per year against 7.4\% for the DAX, with a beta of 0.65, a tracking error of 9.7\% and an information ratio of 0.16. Its CAPM alpha of 4.2\% per year has a $t$-statistic of 2.80.}
\newcommand{\FindSurvivorship}{Using today's members for the whole history raises the equal-weight return from 7.1\% to 11.0\% per year, a bias of 3.8\% per year that would be mistaken for skill.}
\newcommand{\FindSPX}{Measured in euros, the S\&P 500 compounded at 11.5\% per year versus 7.4\% for the DAX, with lower volatility (14.4\% vs 17.7\%). The correlation of monthly returns is 0.69 and the S\&P 500's beta to the DAX is 0.56, so US equities diversified a euro investor's DAX risk rather than replicating it.}
\newcommand{\FindDefensive}{The defensive signals (low volatility and low beta) give long-only portfolios with an average volatility of 13.0\% against 17.7\% for the DAX and an average Sharpe ratio of 0.64, while their long-short spreads average a Sharpe ratio of 0.03: in raw returns their benefit is lower risk rather than a dollar-neutral premium.}
